%
% Se define el tamaño de letra principal:
%
\documentclass[8pt]{extarticle}

%
% Título y autor(es):
%
\title{Currículum Vitae}
\author{von Bergen}

%------------------------- Carga de paquetes --------------------------%
% Se define el tamaño de página y los márgenes:
\usepackage{calc}
\usepackage[a4paper,centering]{geometry}
\setlength{\topmargin}{\topmargin - 0.1\textheight}
\setlength{\textheight}{1.2\textheight}
\setlength{\oddsidemargin}{\oddsidemargin - 0.1\textwidth}
\setlength{\textwidth}{1.2\textwidth}

% Se va a escribir en castellano:
\usepackage[utf8]{inputenc}
\usepackage[\selectedvariety,showlanguages]{babel}
%\usepackage[\selectedvariety , useregional=numeric]{datetime2}
\usepackage[\selectedvariety , \selectedformatdate]{datetime2}
\usepackage{xspace}

\newcommand*{\DTMdisplaydatemmyyyy}[2]{%
	\DTMtwodigits{#2}\DTMsep{yearmonth}\number#1\xspace%
}

% Archivo que define los macros comunes.
\usepackage{cv/common}

% Archivo que define los macros del idioma.
\usepackage{cv/\selectedlanguage /\selectedlanguage}

% Para tener cabecera y pie de página con un estilo personalizado:
\usepackage{fancyhdr}

% Añadir el entorno tabu, permite manejar mejor el texto dentro de tablas.
\usepackage{tabu}
% En el entorno tabular, permite unir filas.
\usepackage{multirow}

% Permite separar el texto en columnas.
\usepackage{multicol}

% Permite modificar el espaciado de las section y subsection.
\usepackage{titlesec}
\titlespacing*{\section}{0pt}{0.5\baselineskip}{0.5\baselineskip}
\titlespacing*{\subsection}{0pt}{0.2\baselineskip}{0.2\baselineskip}
\titleformat*{\subsection}{\large\bfseries}
\titleformat*{\subsection}{\normalsize\bfseries}

\usepackage{amssymb}

% Para enlaces clickeables
\usepackage{xcolor}
\definecolor{DarkBlue}{rgb}{0.0,0.0,0.2} % defino el color azul oscuro
\usepackage[colorlinks=true, linkcolor=DarkBlue, urlcolor=DarkBlue]{hyperref}

%------------------------------ graphicx ----------------------------------
% Para incluir imágenes, el siguiente código carga el paquete graphicx según se está generando un archivo dvi o un pdf (con pdflatex). 
% Para generar dvi, descomente la línea siguiente:
\usepackage{graphicx}
% Para que las imágenes se coloquen en el lugar.
\usepackage{float}
% Para generar pdf (descomente las dos líneas siguientes):
%\usepackage[pdftex]{graphicx}
%\pdfcompresslevel=9

% Todas las imgenes están en el directorio tp-images:
\newcommand{\imgdir}{images}
\graphicspath{{\imgdir/}}

%------------------------- Inicio del documento ---------------------------
\begin{document}
% Hago que el manejo de fin de línea no sea tan estricto.
\sloppy

% Hago que en la cabecera de página se muestre a la derecha la sección, y en el pie, el número de página a la derecha:
\pagestyle{fancy}
\renewcommand{\sectionmark}[1]{\markboth{}{\thesection\ \ #1}}
\lhead{}
\chead{}
\rhead{\rightmark}
\lfoot{}
\cfoot{}
%\rfoot{\thepage\ - \pageref{LastPage}}

% Hago que las páginas se comiencen a contar a partir de aquí:
\setcounter{page}{1}

% Inicio del CV:
\vspace*{-14mm}
\begin{Huge}
\begin{center}
\textbf{\lastname, \name}
\end{center}
\end{Huge}

\leaders\vrule width \textwidth\vskip 0.08pt
\section*{\personalinformation}
\leaders\vrule width \textwidth\vskip 0.08pt

\medskip

\noindent
\begin{minipage}{0.82\linewidth}
\begin{itemize}
\setlength{\itemsep}{0.05cm}
\setlength{\parskip}{0.05cm}
\item \textbf{\phone:} \selectedphonenumbercountrycode \phonenumber
\item \textbf{\mail:} \emailaddress
% \item \textbf{\birthplace:} \birthplaceinfo
\item \textbf{\birthdate:} \birthdateinfo
\item \textbf{\linkedin}
\item \textbf{\GitHub}
\item \textbf{\WebsiteDV}
\end{itemize}
\end{minipage}
\begin{minipage}{0.17\linewidth}
\begin{figure}[H]
	\vspace{-4mm}
	\centering
%	\fbox{\includegraphics[height=0.64\textwidth]{CV_photo}}
	\includegraphics[height=0.64\textwidth]{CV_photo}
	\label{fig: CV_photo}
\end{figure}
\end{minipage}

\medskip

\leaders\vrule width \textwidth\vskip 0.08pt
\section*{\workexperience}
\leaders\vrule width \textwidth\vskip 0.08pt

\medskip

\begin{tabu}{lX[l]c}
% Job VIII DataVision
\textbf{\jobperiodVIII} 
& \textbf{\jobpositionVIII} 
& \multirow{3}{*}{\includegraphics[width=0.12\textwidth]{\jobimageVIII}}\\
	& \jobVIII &\\
 		& \scalebox{0.5}{$\blacksquare$} \textit{\jobinfoVIII} &\\
 		& \scalebox{0.5}{$\blacksquare$} \textit{\jobinfoiVIII} &\\
 		& \scalebox{0.5}{$\blacksquare$} \textit{\jobinfoiiVIII} &\\
 		& \scalebox{0.5}{$\blacksquare$} \textit{\jobinfoiiiVIII} &\\

\\[-0.1pt]

% Job VII Banco Galicia
\textbf{\jobperiodVII} 
& \textbf{\jobpositionVII} 
& \multirow{3}{*}{\includegraphics[width=0.15\textwidth]{\jobimageVII}}\\
	& \jobVII &\\
 		& \scalebox{0.5}{$\blacksquare$} \textit{\jobinfoVII} &\\
 		& \scalebox{0.5}{$\blacksquare$} \textit{\jobinfoiVII} &\\
 		& \scalebox{0.5}{$\blacksquare$} \textit{\jobinfoiiVII} &\\ 

\\[-0.1pt]

% Job VI AB InBev
\textbf{\jobperiodVI} 
& \textbf{\jobpositionVI} 
& \multirow{3}{*}{\includegraphics[width=0.06\textwidth]{\jobimageVI}}\\
	& \jobVI &\\
 		& \scalebox{0.5}{$\blacksquare$} \textit{\jobinfoVI} &\\
 		& \scalebox{0.5}{$\blacksquare$} \textit{\jobinfoiVI} &\\
 		& \scalebox{0.5}{$\blacksquare$} \textit{\jobinfoiiVI} &\\

\\[-0.1pt]

% Job V Mercado Libre
\textbf{\jobperiodV} 
& \textbf{\jobpositionV} 
& \multirow{3}{*}{\includegraphics[width=0.09\textwidth]{\jobimageV}}\\
	& \jobV &\\
 		& \scalebox{0.5}{$\blacksquare$} \textit{\jobinfoV} &\\
 		& \scalebox{0.5}{$\blacksquare$} \textit{\jobinfoiV} &\\
 		& \scalebox{0.5}{$\blacksquare$} \textit{\jobinfoiiV} &\\
 		& \scalebox{0.5}{$\blacksquare$} \textit{\jobinfoiiiV} &\\

\\[0.01pt]

% Job IV Ixpandit
\textbf{\jobperiodIV} 
& \textbf{\jobpositionIV} 
& \multirow{3}{*}{\includegraphics[width=0.14\textwidth]{\jobimageIV}}\\
	& \jobIV &\\
 		& \scalebox{0.5}{$\blacksquare$} \textit{\jobinfoIV} &\\
 		& \scalebox{0.5}{$\blacksquare$} \textit{\jobinfoiIV} &\\
 		& \scalebox{0.5}{$\blacksquare$} \textit{\jobinfoiiIV} &\\
% 			& \texttt{\jobreferenceIV}\\

\\[-0.1pt]
% Job III Inworx
\textbf{\jobperiodIII} 
& \textbf{\jobpositionIII} 
& \multirow{3}{*}{\includegraphics[width=0.14\textwidth]{\jobimageIII}}\\
	& \jobIII &\\
 		& \scalebox{0.5}{$\blacksquare$} \textit{\jobinfoIII} &\\
 		& \scalebox{0.5}{$\blacksquare$} \textit{\jobinfoiIII} &\\
% 			& \texttt{\jobreferenceIII}\\

\\[-0.1pt]

% Job II QBE
\textbf{\jobperiodII} 
& \textbf{\jobpositionII} 
& \multirow{3}{*}{\includegraphics[width=0.11\textwidth]{\jobimageII}}\\
	& \jobII &\\
		& \scalebox{0.5}{$\blacksquare$} \textit{\jobinfoII} &\\
		& \scalebox{0.5}{$\blacksquare$} \textit{\jobinfoiII} &\\
		& \scalebox{0.5}{$\blacksquare$} \textit{\jobinfoiiII} &\\
% 			& \texttt{\jobreferenceII}\\

\\[-0.01pt]

% Job I RBP
\textbf{\jobperiodI} 
& \textbf{\jobpositionI} 
& \multirow{2}{*}{\includegraphics[width=0.07\textwidth]{\jobimageI}}\\
	& \jobI &\\
\end{tabu}

% Inicio Página 2
\newpage

\vspace*{-14mm}

\medskip

\leaders\vrule width \textwidth\vskip 0.08pt
\section*{\computing}
\leaders\vrule width \textwidth\vskip 0.08pt

\medskip

\noindent

% Definimos un comando reutilizable para las secciones
\newcommand{\prolangSection}[3]{%
\begin{minipage}[t]{#1}
\subsection*{\centering{\underline{#2}}}
\smallskip
\begin{itemize}
\setlength{\itemsep}{0.04cm}
\setlength{\parskip}{0.04cm}
#3
\end{itemize}
\end{minipage}
}

% Construcción de las secciones (agrupadas de a tres por línea)
\begin{minipage}[t]{\linewidth}
\prolangSection{0.30\linewidth}{\prolangI}{%
\item Python, R, SQL, PySpark
\item JavaScript
\item HTML, CSS
}
\hfill
\prolangSection{0.30\linewidth}{\prolangII}{%
\item Microsoft Azure
\item Google Cloud Platform (GCP)
\item Amazon Web Services (AWS)
}
\hfill
\prolangSection{0.30\linewidth}{\prolangIII}{%
\item SQL, NoSQL
\item ETL, Data Pipelines
\item BigQuery, Cloud DW
}
\end{minipage}

\vspace{0.3cm}
\noindent

\begin{minipage}[t]{\linewidth}
\prolangSection{0.30\linewidth}{\prolangIV}{%
\item Power BI, Tableau
\item Microsoft Excel
\item Miro
}
\hfill
\prolangSection{0.30\linewidth}{\prolangV}{%
\item Jupyter, PyCharm, VS Code
\item CI/CD for ML pipelines
\item Git/GitHub, Docker, Postman
}
\hfill
\prolangSection{0.30\linewidth}{\prolangVI}{%
\item Jira, Trello
\item Scrum, Kanban
\item CRISP-DM, TDSP
}
\end{minipage}

\medskip

\leaders\vrule width \textwidth\vskip 0.08pt
\section*{\model}
\leaders\vrule width \textwidth\vskip 0.08pt

\medskip

\noindent

% Definimos un comando reutilizable para las secciones de modelos
\newcommand{\modelSection}[3]{%
\begin{minipage}[t]{#1}
\subsection*{\centering{\underline{#2}}}
\smallskip
\begin{itemize}
\setlength{\itemsep}{0.04cm}
\setlength{\parskip}{0.04cm}
#3
\end{itemize}
\end{minipage}
}
% Construcción de las secciones de modelos en Data Science

% Modelos Supervisados
\begin{minipage}[t]{\linewidth}
\modelSection{0.45\linewidth}{\dssegI}{%
\item \dssegImodI
\item \dssegImodII
\item \dssegImodIII
\item \dssegImodIV
\item \dssegImodV
\item \dssegImodVI
}
\hfill
% Modelos IA
\modelSection{0.50\linewidth}{\dssegII}{%
\item \dssegIImodI
\item \dssegIImodII
\item \dssegIImodIII
\item \dssegIImodIV
}
\end{minipage}

\vspace{0.3cm}
\noindent

% Modelos de Deep Learning
\begin{minipage}[t]{\linewidth}
\modelSection{0.45\linewidth}{\dssegVII}{%
\item \dssegVIImodI
\item \dssegVIImodII
}
\hfill
% Modelos de Clustering
\modelSection{0.50\linewidth}{\dssegIV}{%
\item \dssegIVmodI
\item \dssegIVmodII
\item \dssegIVmodIII
}
\end{minipage}

\vspace{0.3cm}
\noindent

% Modelos de Reducción de Dimensionalidad
\begin{minipage}[t]{\linewidth}
\modelSection{0.45\linewidth}{\dssegV}{%
\item \dssegVmodI
\item \dssegVmodII
}
\hfill
% Modelos de Optimización
\modelSection{0.50\linewidth}{\dssegVI}{%
\item \dssegVImodI
\item \dssegVImodII
}
\end{minipage}

\vspace{0.3cm}
\noindent

% Modelos de Series Temporales
\begin{minipage}[t]{\linewidth}
\modelSection{0.45\linewidth}{\dssegIII}{%
\item \dssegIIImodI
\item \dssegIIImodII
}
\hfill
% Modelos de Similitud de Imágenes
\modelSection{0.50\linewidth}{\dssegVIII}{%
\item \dssegVIIImodI
\item \dssegVIIImodII
\item \dssegVIIImodIII
}
\end{minipage}

\medskip

\leaders\vrule width \textwidth\vskip 0.08pt
\section*{\education}
\leaders\vrule width \textwidth\vskip 0.08pt

\medskip
\subsection*{\underline{\textbf{\academic}}}
\begin{tabu}{lX[l]c}
% Titulo Universitario: Actuario
\textbf{\uniperiod} 
& \textbf{\unititle} 
& \multirow{3}{*}{\includegraphics[width=0.04\textwidth]{\uniimageII} \includegraphics[width=0.06\textwidth]{\uniimageI}}\\
 & \uni &\\

% Titulo Secundario: Bachiller
%\textbf{\highschoolperiod} & \textbf{\highschooltitleI} & \multirow{2}{*}{\includegraphics[width=0.03\textwidth]{\highschoolimage}}\\
% & \highschool &\\
\end{tabu}

\subsection*{\underline{\textbf{\courses}}}
\begin{tabu}{l}
\textbf{\courseIII} \\
\textbf{\courseII} \\
\textbf{\courseI} 
\end{tabu}

\medskip

\leaders\vrule width \textwidth\vskip 0.08pt
\section*{\languages}
\leaders\vrule width \textwidth\vskip 0.08pt

\medskip

\begin{tabu} to \textwidth {@{}l p{\dimexpr\textwidth-6cm}@{}}
\textbf{\langI}
&
\raggedleft
\multirow{3}{*}{\includegraphics[width=0.1\textwidth]{\languageimage}}
\\
\langexamiI
&
\\
&
\\
\end{tabu}

%\begin{tabu}{l r}
% Ingles 
%\textbf{\langexamidateI} 
%\textbf{\langI} 
%& \multirow{3}{*}{\includegraphics[width=0.12\textwidth]{\languageimage}} \\
%\langexamiI
%\end{tabu}

\medskip

\leaders\vrule width \textwidth\vskip 0.08pt
\section*{\projects}
\leaders\vrule width \textwidth\vskip 0.08pt

\medskip

\begin{tabu} to \textwidth {@{}l p{\dimexpr\textwidth-13cm}@{}}

% Prendiendo empatia
\textbf{\projectsII}
&
\raggedleft
\multirow{3}{*}{\includegraphics[width=0.08\textwidth]{\projectsimageII}}
\\
\scalebox{0.5}{$\blacksquare$} \textit{\projectsinfoII}
&
\\
\scalebox{0.5}{$\blacksquare$} \textit{\projectsinfoiII}
&
\\[0.5pt]

% Ciclo Mestiza
\textbf{\projectsI}
&
\raggedleft
\multirow{3}{*}{\includegraphics[width=0.07\textwidth]{\projectsimageI}}
\\
\scalebox{0.5}{$\blacksquare$} \textit{\projectsinfoI}
&
\\

\end{tabu}
\medskip


%\begin{flushleft}
%\begin{tabular}{lp{0.79\textwidth}}
%\textbf{\addinfoperiodII} & \textbf{\addinfoII}\\
%{\includegraphics[width=0.12\textwidth]{\addinfoimageI}}\\
% & \addinfoinfoII\\
% & \scalebox{0.5}{$\blacksquare$} \textit{\addinfodescriptionII}\\

%\textbf{\addinfoperiodI} & \textbf{\addinfoI}\\
% & \addinfoinfoI\\
% & \scalebox{0.5}{$\blacksquare$} \textit{\addinfodescriptionI}\\

%\end{tabular}
%\end{flushleft}


%\begin{minipage}[t]{1\linewidth}
%	\smallskip
%	\begin{minipage}[t]{1\linewidth}
%		\begin{itemize}
%			\setlength{\itemsep}{0.04cm}
%			\setlength{\parskip}{0.04cm}
%			\item  I agree to the processing of my personal data by 90mph Consulting Lukasz Bielawski based in Legionowo 05-120, Mieszka 1/25 street and Business Partners for the needs of this recruitment process and any future recruitments in accordance with Personal Data Protection Act of 10 May 2018 (Journal of Laws 2018, item 1000) and in agreement with Regulation (EU) 2016/679 of the European Parliament and of the Council of 27 April 2016 on the protection of natural persons with regard to the processing of personal data and on the free movement of such data, and repealing Directive 95/46/EC (General Data Protection Regulation). At the same time I declare that I have been informed that: - The administrator of my personal data is 90mph Consulting Lukasz Bielawski; - I have the right to access my data, as well as rectify, delete or limit processing. Personal data may be transferred to authorized entities, e.g. potential employers on the basis of legal provisions.
%I hereby consent to my personal data being processed by 90mph Consulting Lukasz Bielawski, Mieszka I/25 street, 05-120 Legionowo for the purpose of considering my application for the vacancy Expert Credit Risk Modelling, Expert Credit Model Validationt and else vacancies in ING Tech Poland and ING Business Shared Services B.V. spolka z ograniczona odpowiedzialnoscia.
%		\end{itemize}
%	\end{minipage}
%\end{minipage}
%\vrule width 0.04pt

\end{document}
